\documentclass[a4paper,9pt]{extarticle}

% Packages
\usepackage[utf8]{inputenc} % For input encoding
\usepackage{geometry} % For page margins
\geometry{a4paper, margin=0.9in} % Set paper size and margins
\usepackage{titlesec} % For section title formatting
\usepackage{enumitem} % For itemized list formatting
\usepackage{color}
\usepackage{hyperref} % For hyperlinks
\hypersetup{
  hidelinks,
  colorlinks=true,
  linkcolor=blue,
  citecolor=blue,
  urlcolor=blue
}

% Formatting
\setlist{noitemsep} % Removes item separation
\titleformat{\section}{\large\bfseries}{\thesection}{1em}{}[\titlerule] % Section title format
\titlespacing*{\section}{0pt}{\baselineskip}{\baselineskip} % Section title spacing

% Begin document
\begin{document}

% Disable page numbers
\pagestyle{empty}

% Header
\begin{center}
\textbf{\Large Hanjun Luo}\\[3pt]
New York City, US \,|\, \href{mailto:hl6266@nyu.edu}{hl6266@nyu.edu} \,|\, 
\href{https://scholar.google.com/citations?user=RcWubw8AAAAJ&hl=en}{Google Scholar} \,|\, 
\href{https://github.com/Astarojth}{GitHub} \,|\, 
\href{https://astarojth.github.io/}{Homepage}
\end{center}

%------------------------
% Education Section
\section*{EDUCATION}
\noindent
\textbf{Ph.D.~in Computer Science} \\
\textbf{New York University} \\
Expected completion: 2030\\

\noindent
\textbf{Bachelor of Engineering in Computer Engineering} \\
\textbf{Zhejiang University} \\
Year of completion: 2025\\

\noindent
\textbf{Bachelor of Science in Computer Engineering} \\
\textbf{University of Illinois Urbana-Champaign} \\
Year of completion: 2025

%------------------------
% Research Interests
\section*{RESEARCH INTERESTS}
\vspace{-0.3\baselineskip}
\begin{itemize}[topsep=3pt,itemsep=1pt]
  \item Evaluation of LLMs, MLLMs, and LLM agents
  \item LLM/agent-as-a-judge methodology
  \item Human-agent collaboration
  \item Safety of LLM agents
  \item Bias and trustworthiness in generative models
  \item LLM-based multilingual information extraction
\end{itemize}

%------------------------
% Publications Section
\section*{PUBLICATIONS}
\noindent
\textbf{Luo, H.}, Huang, Z., Chung, S., Wang, Y., Jin, Y., Li, J., Li, J., Li, X., \& Salam, H. (2026). AtelierEval: Agentic Evaluation of Humans \& LLMs as Text-to-Image Prompters. \textit{International Conference on Machine Learning (ICML 2026)}.\\[4pt]
\noindent
\textbf{Luo, H.}, Ni, C., Wen, J., Huang, Z., Wang, Y., Liao, B., Chung, S., Jin, Y., Li, J., Li, X., Xu, W., Wang, X., \& Salam, H. (2026). CentaurEval: Benchmarking Human-in-the-Loop Value in Agentic Coding. \textit{International Conference on Machine Learning (ICML 2026)}.\\[4pt]
\noindent
\textbf{Luo, H.}, Dai, S., Ni, C., Li, X., Zhang, G., Wang, K., Liu, T., \& Salam, H. (2025). AgentAuditor: Human-level safety and security evaluation for LLM agents. \textit{Conference on Neural Information Processing Systems (NeurIPS 2025), poster}.\\[4pt]
\noindent
\textbf{Luo, H.}, Jin, Y., Li, X., Liu, X., Chen, R., Shang, T., Wang, K., Wen, Q., \& Liu, Z. (2025). DynamicNER: A dynamic, multilingual, and fine-grained dataset for LLM-based named entity recognition. \textit{Conference on Empirical Methods in Natural Language Processing (EMNLP 2025), main conference}.\\[4pt]
\noindent
\textbf{Luo, H.}, Huang, Z., Huang, H., Deng, Z., Chen, R., Li, X., Liu, Z., \& Salam, H. (2026). BiasIG: Benchmarking multi-dimensional social biases in text-to-image models. \textit{International Joint Conference on Neural Networks (IJCNN 2026)}.\\[4pt]
\noindent
\textbf{Luo, H.}, Deng, Z., Chen, R., \& Liu, Z. (2024). FAIntbench: A holistic and precise benchmark for bias evaluation in text-to-image models. \textit{ICML 2024 Workshop on Data-centric Machine Learning Research (DMLR)}.\\[4pt]
\noindent
Li, K., Shen, C., Liu, Y., Han, J., Zheng, K., Zou, X., Wang, Z., Du, X., Zhang, S., \textbf{Luo, H.}, \textit{et al.} (2026). AudioTrust: Benchmarking the multifaceted trustworthiness of audio large language models. \textit{International Conference on Learning Representations (ICLR 2026), poster}.\\[4pt]
\noindent
Li, J., Chen, Z., \textbf{Luo, H.}, \& Salam, H. (2026). PrefIx: Understand and adapt to user preference in human-agent interaction. \textit{Annual Meeting of the Association for Computational Linguistics (ACL 2026 Findings)}.\\[4pt]
\noindent
Jin, Y., Du, X., \textbf{Luo, H.}, Wang, Z., Hu, H., Wang, X., \& Li, X. (2026). AudioStealer: Extracting audio prompts via Shapley value-guided query search. \textit{Annual Meeting of the Association for Computational Linguistics (ACL 2026 Findings)}.\\[4pt]
\noindent
Sun, M., Zhao, Z., Chai, W., \textbf{Luo, H.}, Cao, S., Zhang, Y., Hwang, J.-N., \& Wang, G. (2024). UniAP: Towards universal animal perception in vision via few-shot learning. \textit{AAAI Conference on Artificial Intelligence (AAAI 2024)}.\\[4pt]

%------------------------
% Research Experience Section 
\section*{RESEARCH EXPERIENCE}
\noindent
\textbf{LLM/MLLM/Agent Evaluation}\\
-- AtelierEval (ICML 2026): Build a benchmark with 360 expert-crafted tasks to evaluate text-to-image prompting proficiency for both MLLMs and human prompters.\\
-- CentaurEval (ICML 2026): Build a coding benchmark for both LLM agents and human developers with 45 collaboration-necessary templates and 450 instantiated tasks, showing low pass rates for LLM-only (0.67\%) and human-only (18.89\%) settings but a clear collaborative gain to 31.11\%.\\
-- AgentAuditor (NeurIPS 2025): Build ASSEBench, an LLM-as-a-Judge benchmark with 2,293 annotated interaction records covering 15 risk types across 29 application scenarios for end-to-end agent safety/security evaluation.\\
-- DynamicNER (EMNLP 2025): Build a dynamic and multilingual benchmark spanning 8 languages, 155 entity types, and 3 granularity levels to evaluate robustness and generalization of LLM-based NER.\\
-- PrefIx (ACL Findings 2026): Introduce a configurable benchmark with Interaction-as-a-Tool (IaaT), 31 preference settings, and 14 attributes to evaluate LLM-agent preference adaptation beyond task accuracy.\\
-- AudioTrust (ICLR 2026): Construct a multifaceted benchmark spanning 6 trustworthiness dimensions, 26 sub-tasks, and 4,420+ audio samples, with evaluation across 14 SOTA ALLMs under 18 experimental configurations.\\[-2pt]

\noindent
\textbf{LLM/Agent-as-a-Judge}\\
-- AgentAuditor (NeurIPS 2025): Introduce a novel Agent-as-a-Judge method with memory-augmented reasoning for agent safety and security evaluation.\\
-- BiasIG (IJCNN 2026): Construct a fine-tuned MLLM-based judge for bias in text-to-image models.\\
-- PrefIx (ACL Findings 2026): Propose a composite LLM-as-a-Judge protocol over seven UX dimensions for scalable preference-alignment scoring in agent interactions.\\[-2pt]


\noindent
\textbf{LLM Agent Safety}\\
-- AgentAuditor (NeurIPS 2025): Introduce both an Agent-as-a-Judge method and a benchmark for agent safety and security evaluation.\\
-- A Comprehensive Survey in LLM(-agent) Full Stack Safety: Lead the benchmarking component and systematize LLM and agent safety risks and mitigation across data, training, and deployment.\\[-2pt]

\noindent
\textbf{Human--Agent Collaboration}\\
-- CentaurEval (ICML 2026): Design collaboration-necessary tasks to quantify human-in-the-loop value in coding and verify them through human studies and LLM benchmarking.\\
-- PrefIx (ACL Findings 2026): Propose a preference-aware LLM-agent collaboration method that infers implicit user preferences and adapts confirmation style and pacing in multi-turn workflows.\\[-2pt]

\noindent
\textbf{Bias Evaluation and Debiasing in Text-to-Image Models}\\
-- BiasIG (IJCNN 2026): Construct a multi-dimensional social-bias benchmark with 47,040 prompts across 4 dimensions and benchmark 8 T2I models with 3 debiasing methods.\\
-- FAIntbench (ICML-DMLR 2024): Build a holistic benchmark for fine-grained text-to-image bias evaluation.\\
-- VersusDebias: Propose a zero-shot debiasing framework and analyze trade-offs.\\[-2pt]

\noindent
\textbf{Audio Model Trustworthiness}\\
-- AudioTrust (ICLR 2026): Based on the benchmark, identify key findings on audio-specific vulnerabilities from non-semantic cues (e.g., accent, timbre, background noise) across 14 ALLMs and 18 evaluation settings.\\
-- AudioStealer (ACL Findings 2026): Design a two-stage Shapley-guided black-box prompt extraction attack. Introduce Prompt2Music, a prompt-stealing benchmark with 5,000 prompt-audio pairs generated from 4 text-to-music models.\\[2pt]

\noindent
\textbf{LLM-based Multilingual Information Extraction}\\
-- DynamicNER (EMNLP 2025): Develop a realistic multilingual setting for NER, covering evolving entities and fine-grained categories.

\section*{VISITING EXPERIENCE}
\noindent
\textbf{Tsinghua University} \\
2026.5--2026.8 \\
\textit{Visiting Researcher}, Prof. Yipeng Dong

% Teaching Experience Section 
\section*{TEACHING EXPERIENCE}
\noindent
\textbf{CS101: Intro to Computing} \\
2024.9--2024.12 \\
\textit{Teaching Assistant}, Zhejiang University, Prof. Wee Liat Ong

%------------------------
% Professional Skills Section 
\section*{ADDITIONAL SKILLS}
\vspace{-0.3\baselineskip}
\begin{itemize}[topsep=3pt,itemsep=1pt]
  \item \textbf{Languages}: Chinese (native), English (fluent).
  \item \textbf{Programming}: Python, SQL, LaTeX.
  \item \textbf{Photography}: Signed contributor for Visual China Group (VCG); experienced in landscape photography.
\end{itemize}

\end{document}
